\section{Related Work}
\label{sec:related}

\paragraph{Sampling-time text watermarks.}
Green-list watermarking partitions the vocabulary at each step with a keyed hash of recent context and biases sampling toward the green set; detection counts green tokens against a null \citep{kgw}. SynthID-Text generalises the same idea with a tournament sampling procedure designed for production deployment and for detection from text alone \citep{synthidtext}. Exponential and Gumbel schemes key the sampler's noise and admit exact tail tests \citep{kuditipudi}. Distortion-free schemes sample from the original distribution while retaining a detectable coupling \citep{christ}. Our study targets the SynthID-Text class because it is deployed, because an open reference implementation exists, and because the detection test is well specified; we expect the qualitative ranking of tactics to transfer to other sampling-bias schemes, and we treat that transfer as a question for future work.

\paragraph{Removal attacks.}
Paraphrasing is the canonical attack. \citet{kgwreliability} show that machine paraphrase reduces detection at short lengths and that the mark recovers as the passage grows, and that human rewriting is stronger than machine rewriting. DIPPER, an 11B discourse-level paraphraser, evades several detectors including green-list watermarking while preserving meaning closely \citep{dipper}. Token-substitution attacks replace a fraction of words with synonyms and were used as the robustness probe in \citet{threebricks}. A more targeted line masks high-self-information tokens and infills them, which is the masked-LM infilling row of our catalogue \citep{sira}. Regeneration with a different model replaces the token stream wholesale. The impossibility result of \citet{sand} gives a quality-oracle random walk that removes any watermark that occupies a small share of the high-quality output space, and \citet{sadasivan} argue more broadly that detection is fragile under paraphrase. \citet{jovanovic} study the dual problem of stealing a scheme and forging marks, show that stealing also boosts scrubbing, and bound what a detector can be trusted to report. An independent evaluation applies the attack to SynthID-Text and reports that an off-the-shelf paraphraser clears the mark in over 90\% of roughly 1,000-token samples at a \(10^{-3}\) false-positive rate \citep{probingsynthid}.

\paragraph{Back-translation and cross-lingual removal.}
Round-trip translation through a pivot language perturbs token choice while preserving propositional content, which makes it an attractive removal primitive. \citet{cwra} formalise the cross-lingual removal attack, show that watermark strength is not consistent across languages, and report a large drop in detection after a pivot round trip. MarkMyWords includes translation among its tamper-resistance perturbations and reports that long generations often remain detectable \citep{markmywords}. The two results differ in scheme, length, and pivot, which is one reason a joint evaluation under a common operating point is useful.

\paragraph{Benchmarks and tooling.}
MarkMyWords unifies several schemes behind one interface and reports quality, detection length, and tamper resistance \citep{markmywords}. MarkLLM provides reference implementations of the KGW and SynthID scheme families together with generation, detection, and visualisation utilities \citep{markllm}. We build on MarkLLM for the detector and keep our harness independent of any single scheme so that the catalogue and the search can be reused.

\paragraph{Evaluation of preservation.}
Surface overlap such as BLEU is a poor proxy for meaning under paraphrase. We therefore combine embedding-based similarity \citep{bertscore}, learned reference-based quality \citep{bleurt}, learned translation quality \citep{comet}, and distributional closeness \citep{mauve}. Section~\ref{sec:metrics} states how the signals are aggregated.
